%% kapa: exam=1 round=31 year=2020 period=1 form=A subject=real_estate subject_ko=부동산학원론 seq=38 number=118
%% source: Q-Net 감정평가사 기출문제 (cst00309) artlSeq=5207219 file=extracted/1008/5207219/A형/A형/2020년도 제31회 감정평가사 1차 시험_1교시_A형.pdf
%% conversion: kordoc 4.24.51 -> tools/gen1.py (auto); review=proofread-sonnet; flags=none
\begin{question}{118}
부동산 거래신고 등에 관한 법률상 옳지 않은 것은? (단, 주어진 조건에 한함)
\begin{choices}
  \item 거래당사자 중 일방이 지방자치단체인 경우에는 지방자치단체가 신고를 하여야 한다.
  \item 공동으로 중개한 경우에는 해당 개업공인중개사가 공동으로 신고하여야 하며, 일방이 신고를 거부한 경우에는 단독으로 신고할 수 있다.
  \item 거래당사자는 그 실제 거래가격 등을 거래계약의 체결일부터 30일 이내에 공동으로 신고해야 한다.
  \item 누구든지 개업공인중개사에게 부동산 거래의 신고를 하지 아니하게 하거나 거짓으로 신고하도록 요구하는 행위를 하여서는 아니 된다.
  \item 거래당사자가 부동산의 거래신고를 한 후 해당 거래계약이 취소된 경우에는 취소가 확정된 날부터 60일 이내에 해당 신고관청에 공동으로 신고하여야 한다.
\end{choices}
\end{question}
\answer{5}
